\section{Discussion}

The central methodological result is not merely that one restricted search happens to be exact. It is that feasibility, exact cost, structural explanation, and finite row verification can be assigned separate authorities. The separation makes a counterexample diagnostic. The first adverse row identified a missing borrow geometry. The hostile panel then showed that the first repair was also incomplete. In both cases the support-two evaluator remained exact because its certificate did not depend on the explanation family.

This design also prevents a common inversion of evidence. The 9,547 zero-error comparison events are strong implementation evidence, but they cannot prove behavior for arbitrary qubit count. The proof reduces the all-size claim to a combinatorial subset lemma and a finite local inequality. Conversely, the proof does not certify that a particular executable, serialization, or table was produced correctly; those are checked by independent witnesses and arithmetic consistency tests.

The explanation theorem adds value beyond exact cost. A family label exposes whether the minimizer uses independent weight-one anchors, an out-of-support borrow, or a two-anchor hybrid. Such labels can guide compiler diagnostics and future simplification, but they should not be confused with physical mechanisms. They classify feasible configurations within the abstract grammar of Eq.~\eqref{eq:cost}.

Negative proof routes are scientifically useful here. A replacement-orbit argument failed because its ten candidate images returned to the same qubit pair and provided no strict descent. A separate single-pinner normalization survived a complete 133,349,376-configuration product domain, yet its proposed global chain composition remained unproved. A later exact three-tag configuration refuted the premise that every chain can be represented on two coordinates. These failures do not weaken the successful per-block-permutation proof; they explain why apparently simpler globalizations were not used. Appendix~\ref{app:explanation-audit} records their logical status.

The objective counterexample is equally important. Equation~\eqref{eq:reweighted} shows that ``support two is enough'' is not a property of the grammar alone. It results from the exact balance between frame refund and local restoration increase in the unit objective. Reweighting can make a larger support profitable. Reporting this boundary is preferable to suggesting robustness that has not been proved.

No canonical speed comparison is claimed. Restricted enumeration may be substantially cheaper than unrestricted optimization in some regimes, but runtime depends on implementation choices, caching, symmetry reductions, and hardware. Correctness and speed should be benchmarked separately on the final public implementation.
